% =====================================================================
% SEÇÃO: Arquitetura Transformer (para integrar ao Capítulo 2)
% Criada em rodada separada para revisão isolada.
% Integrar posteriormente em 02_capitulo/fundamentacao_teorica.tex,
% idealmente ANTES da seção "Consumo Energético em Redes Neurais".
% =====================================================================

\section{Arquitetura Transformer}\label{sec:transformer}

A arquitetura \textit{Transformer}, proposta por \citeonline{Vaswani2017} no trabalho
\textit{Attention is All You Need}, marcou uma mudança profunda na forma como modelos
de aprendizado profundo processam sequências de texto. Antes dessa proposta, o campo
era dominado por arquiteturas que processavam a sequência elemento por elemento, palavra
por palavra. O \textit{Transformer} rompeu com essa abordagem ao introduzir um mecanismo
capaz de considerar toda a sequência de uma só vez. Esta seção apresenta, de forma
progressiva, os conceitos necessários para compreender a arquitetura utilizada neste
trabalho: as limitações das abordagens anteriores, o mecanismo de atenção, a operação
de autoatenção e sua versão multi-cabeça, a estrutura de cada bloco da rede e,
finalmente, o modelo GPT-2, que é o objeto central dos experimentos de poda.

% COLOCAR FIGURA: diagrama simplificado comparando o processamento sequencial de
% uma RNN (palavras processadas uma a uma, da esquerda para a direita) com o
% processamento paralelo do Transformer (todas as palavras conectadas entre si
% simultaneamente). Elaboração própria.

\subsection{Limitações das Arquiteturas Recorrentes}\label{subsec:limitacoes_rnn}

Até a chegada do \textit{Transformer}, o processamento de texto era dominado pelas redes
neurais recorrentes (\textit{Recurrent Neural Networks}, RNNs), popularizadas ao longo da
década de 1990. Diferentemente de arquiteturas puramente diretas, as RNNs processam a
sequência palavra por palavra e reaproveitam a saída de cada passo como entrada do passo
seguinte, criando uma forma de ``memória'': a cada nova palavra, um estado oculto
(\textit{hidden state}) é atualizado com base na palavra atual e no estado anterior, de
modo que, ao final, ele deve resumir as informações relevantes de todo o texto.

Essa memória, porém, falha em sequências longas. Considere o trecho ``A Inglaterra é minha
terra natal. Passei toda a minha vida lá. Mudei-me para a Espanha há dois dias. Eu falo
apenas uma língua, que é o \_\_\_''. Para completar a lacuna corretamente com ``inglês'', o
modelo precisa recuperar a palavra ``Inglaterra'', que aparece no início do trecho. O
obstáculo é o chamado gradiente desaparecente (\textit{vanishing gradient}): durante o
treinamento via \textit{backpropagation}, os gradientes são multiplicados repetidamente
pelas matrizes de peso da rede e, quando esses valores são pequenos, encolhem até se
aproximarem de zero \cite{Bengio1994}. Com isso, a influência de ``Inglaterra'' torna-se
insignificante ao chegar à lacuna, e o modelo perde o contexto necessário --- limitação
que se manifesta como incapacidade de lidar com \textbf{dependências de longo alcance}.

Há ainda uma segunda limitação, de ordem computacional: como o estado oculto no instante
$t$ só pode ser calculado após o do instante $t-1$, o processamento é estritamente
sequencial e mal aproveita o paralelismo das unidades de processamento gráfico (GPUs)
modernas. Essas duas dificuldades impulsionaram uma evolução cronológica de arquiteturas.
As redes LSTM (\textit{Long Short-Term Memory}), propostas em 1997, atacaram o gradiente
desaparecente por meio de portas (\textit{gates}) que regulam o que deve ser mantido ou
esquecido, permitindo reter por mais tempo informações como ``Inglaterra''
\cite{Hochreiter1997}. Ainda assim, as LSTMs continuam processando a frase de forma
sequencial, o que as mantém lentas e difíceis de paralelizar em grandes volumes de dados
--- gargalo que só seria superado pelo mecanismo de atenção.

% SUGESTÃO DE FIGURA: diagrama mostrando uma RNN processando a frase de exemplo,
% com setas indicando o fluxo do estado oculto e destacando a "distância" entre
% "Inglaterra" e a lacuna final. Elaboração própria.

\subsection{Mecanismo de Atenção}\label{subsec:atencao}

O passo seguinte dessa evolução foi o mecanismo de atenção (\textit{attention
mechanism}), introduzido por \citeonline{Bahdanau2015} no contexto da tradução automática.
Em vez de comprimir toda a frase em um único estado oculto, a atenção permite que o modelo
``olhe'' de volta e foque seletivamente em partes específicas da entrada. No exemplo
anterior, ela pode consultar diretamente o termo ``Inglaterra'', independentemente de quão
distante ele esteja da lacuna, eliminando o risco de o gradiente desaparecer pelo caminho.

Para operacionalizar essa ideia, o mecanismo utiliza três conceitos: a consulta
(\textit{query}), a chave (\textit{key}) e o valor (\textit{value}). Uma analogia útil é a
de uma biblioteca: cada livro tem uma etiqueta na lombada (a \textit{chave}) que descreve
seu conteúdo; quando um leitor chega com uma pergunta (a \textit{consulta}), ela é comparada
com cada etiqueta para decidir quais livros são mais relevantes; o conteúdo desses livros
(os \textit{valores}) é então combinado, com os mais relevantes contribuindo mais para a
resposta. No \textit{Transformer}, esse processo ocorre com números: a relevância entre uma
consulta e cada chave é calculada matematicamente, e os valores são somados com pesos
proporcionais a essa relevância.

A inovação central de \citeonline{Vaswani2017}, em 2017, foi construir uma arquitetura
inteira baseada exclusivamente nesse mecanismo, dispensando por completo a recorrência.
Com isso, todas as posições de uma sequência se comunicam diretamente em uma única
operação, sem precisar ``passar'' por palavras intermediárias. Isso resolve as duas
limitações das RNNs de uma só vez: as dependências de longo alcance tornam-se tão fáceis
de capturar quanto as de curto alcance, e o cálculo passa a ser paralelizável, tornando
as arquiteturas baseadas em RNN e LSTM obsoletas para muitas tarefas modernas.

\subsection{Autoatenção e Atenção Multi-Cabeça}\label{subsec:autoatencao}

A variante do mecanismo de atenção utilizada no \textit{Transformer} é chamada de
autoatenção (\textit{self-attention}). O prefixo ``auto'' indica que a consulta, a chave
e o valor são todos derivados da mesma sequência de entrada. Ou seja,
ao processar uma frase, cada palavra pergunta a si mesma: ``quais outras palavras desta
frase são relevantes para minha interpretação?''. Exemplo: ``O gato que estava
no telhado dormiu''. Ao processar a palavra ``dormiu'', o mecanismo de autoatenção
calcula um peso para cada palavra da frase. A palavra ``gato'' recebe peso alto, pois é
o sujeito do verbo; palavras como ``telhado'' ou ``no'' recebem pesos baixos, pois são
menos relevantes para determinar quem dormiu.

Na prática, esses pesos são calculados da seguinte forma, a partir de cada representação
de entrada, três matrizes de pesos aprendíveis durante o treinamento transformam cada
representação em três vetores: o vetor de consulta $Q$ (\textit{query}), o vetor de
chave $K$ (\textit{key}) e o vetor de valor $V$ (\textit{value}). O grau de compatibilidade
entre uma consulta e cada chave é calculado pelo produto escalar entre eles. Esses valores
brutos de compatibilidade passam pela função \textit{softmax}, que os converte em uma
distribuição de probabilidades, isto é, um conjunto de números entre 0 e 1 que somam
exatamente 1. Esses pesos normalizados determinam com que intensidade cada valor $V$
contribui para o resultado final, que é obtido pela soma ponderada dos valores.

Essa operação completa é denominada \textit{Scaled Dot-Product Attention} (atenção por
produto escalar escalado) e é definida formalmente por \cite{Vaswani2017}:

\[
    \text{Attention}(Q, K, V) = \text{softmax}\!\left(\frac{QK^\top}{\sqrt{d_k}}\right) V
\]

Em termos intuitivos, o produto $QK^\top$ mede a compatibilidade entre todas as consultas e
todas as chaves de uma só vez; a divisão por $\sqrt{d_k}$ (a dimensão das chaves) é um fator
de escala que estabiliza o treinamento; a \textit{softmax} converte essas pontuações em
pesos que somam 1; e a multiplicação por $V$ produz, para cada posição, uma média ponderada
dos valores, isto é, sua representação já contextualizada pelo restante da sequência.

Uma única operação de atenção, porém, captura apenas um tipo de relação entre as palavras
de cada vez. Para superar essa limitação, \citeonline{Vaswani2017} propõem a atenção
multi-cabeça (\textit{multi-head attention}). A analogia é a seguinte: se a atenção
simples é como um único leitor relendo a frase em busca de relações gramaticais, a atenção
multi-cabeça é como $h$ leitores independentes, cada um procurando um tipo diferente de
relação. Um leitor pode focar em concordância verbal (``gato''/``dormiu''), outro em
relações de posse (``telhado do vizinho''), outro em referências pronominais, e assim por
diante. O GPT-2 \textit{small}, por exemplo, utiliza 12 cabeças de atenção por bloco.

Na prática, cada uma das $h$ cabeças executa a operação \textit{Scaled Dot-Product
Attention} de forma independente, com suas próprias matrizes de projeção aprendíveis para
$Q$, $K$ e $V$. As saídas das $h$ cabeças são depois concatenadas e projetadas por uma
matriz linear aprendível, produzindo a representação final de cada posição. Desse modo,
o modelo pode simultaneamente capturar relações em diferentes subespaços semânticos e
sintáticos, sem custo adicional expressivo em relação à atenção simples.

Sob a perspectiva da compressão de modelos, as cabeças de atenção constituem unidades
podáveis naturais para a poda estruturada: estudos da literatura demonstram que nem todas
as cabeças são igualmente úteis, e que algumas podem ser removidas inteiramente sem
degradação significativa da qualidade das predições \cite{Michel2019}. 
% Essa propriedade
% é explorada diretamente nos experimentos deste trabalho.

% SUGESTÃO DE FIGURA: diagrama comparativo entre Scaled Dot-Product Attention (esquerda,
% mostrando Q, K, V, o produto escalar, a softmax e a ponderação de V) e Multi-Head
% Attention (direita, mostrando h cabeças paralelas com concatenação e projeção final),
% adaptado de Vaswani et al. (2017), Fig. 2.

\subsection{Estrutura do Bloco Transformer}\label{subsec:bloco_transformer}

Para entender como o \textit{Transformer} processa informação em profundidade, é útil
pensar em cada bloco como uma ``estação de trabalho'' em uma linha de montagem. Na
primeira bancada, as palavras trocam informações entre si por meio da atenção
multi-cabeça: cada palavra atualiza sua representação consultando todas as demais. Na
segunda bancada, cada palavra refina individualmente sua própria representação, de forma
independente das outras, por meio de uma pequena rede totalmente conectada.
Esse par de operações se repete em cada bloco, e os blocos são empilhados, fazendo com
que as representações das palavras se tornem progressivamente mais ricas e abstratas a
cada andar.

O primeiro bloco é a
camada de atenção multi-cabeça, já apresentada na seção anterior, responsável por
integrar informação contextual de toda a sequência. O segundo é uma rede
\textit{feedforward} totalmente conectada --- um MLP (\textit{multilayer perceptron})
aplicado de forma idêntica e independente a cada posição. Esse MLP é composto por duas
transformações lineares com uma função de ativação não linear entre elas: \modificado{a função GELU
(\textit{Gaussian Error Linear Unit}) \cite{Hendrycks2016} nos modelos GPT, ou ReLU em outras variantes.
\adicionadomarcelo{Diferentemente da ReLU \cite{Nair2010}, que simplesmente zera valores negativos, a
GELU faz uma transição suave em torno da origem. Na prática, essa suavidade evita o problema dos
``neurônios mortos'' (\textit{dying ReLU}), em que unidades deixam de aprender por receberem gradiente
nulo, e por isso \citeonline{Radford2019} a adotaram como ativação padrão na família GPT.}} A
dimensão intermediária desse MLP é tipicamente quatro vezes maior do que a dimensão
do modelo --- no GPT-2 \textit{small}, a dimensão do modelo é 768 e a dimensão
intermediária é 3.072 --- o que confere ao bloco grande capacidade de representação antes
de comprimir novamente a informação.

% Cada um dos dois sub-módulos é envolvido por dois mecanismos que facilitam o treinamento
% de modelos profundos. O primeiro é a conexão residual (\textit{residual connection}),
% idêntica em conceito à utilizada nas redes ResNet \cite{He2016}: a entrada do
% sub-módulo é somada diretamente à sua saída, criando um ``atalho'' que permite ao gradiente
% fluir mais facilmente para camadas anteriores durante o \textit{backpropagation}. O leitor
% familiarizado com arquiteturas convolucionais reconhecerá esse mecanismo. O segundo é a
% normalização de camada (\textit{layer normalization}), que estabiliza as ativações ao
% longo do treinamento, normalizando-as ao longo da dimensão das representações em vez da
% dimensão do lote, tornando o treinamento mais estável e menos sensível à taxa de
% aprendizado.

% A arquitetura completa de um modelo baseado em \textit{Transformer} consiste em um número
% $L$ de blocos empilhados. O GPT-2 \textit{small} possui $L = 12$ blocos. Cada bloco
% adicional aumenta a capacidade do modelo de aprender representações mais abstratas, da
% mesma forma que camadas mais profundas em uma rede convolucional capturam características
% visuais progressivamente mais complexas. Essa profundidade, porém, vem com custo
% computacional e energético proporcional: cada bloco extra adiciona parâmetros e operações
% ao caminho de cada inferência.

% Do ponto de vista da compressão, essa organização modular expõe múltiplos alvos para a
% poda estruturada: os neurônios intermediários da camada MLP (que podem ser removidos
% reduzindo a dimensão intermediária), as cabeças de atenção (que podem ser eliminadas
% inteiramente) e, em cenários de compressão mais agressiva, blocos \textit{Transformer}
% inteiros. Essa variedade de alvos torna a poda estruturada em modelos \textit{Transformer}
% mais rica em granularidade do que em redes convolucionais como a ResNet-50, onde os
% alvos típicos são filtros e canais de convolução \cite{Li2017}.

% SUGESTÃO DE FIGURA: diagrama de um único bloco Transformer (decoder-only), mostrando
% a camada de atenção multi-cabeça, a camada MLP feedforward, as conexões residuais e
% a normalização de camada, com as dimensões do GPT-2 small anotadas (768, 3072, 12 cabeças).
% Elaboração própria.

\subsection{Modelos de Linguagem de Grande Escala e o GPT-2}\label{subsec:llm_gpt2}

Modelos de linguagem de grande escala (\textit{Large Language Models}, LLMs) são redes
neurais profundas, tipicamente baseadas na arquitetura \textit{Transformer}, treinadas em
grandes volumes de texto com o objetivo de aprender a prever palavras. A tarefa de
pré-treinamento é simples de enunciar: dado um trecho de texto, o modelo deve prever
qual token vem a seguir. Um \textit{token} é a unidade mínima de texto que o modelo
processa --- pode ser uma palavra, parte de uma palavra, um sinal de pontuação ou um
espaço. Dada a sequência ``O aluno abriu o'', um LLM bem treinado atribuiria alta
probabilidade a tokens como ``livro'', ``caderno'' ou ``computador'', e baixa probabilidade
a tokens como ``foguete'' ou ``plutônio''. Repetida bilhões de vezes sobre textos variados,
essa tarefa aparentemente simples força o modelo a aprender gramática, fatos sobre o mundo,
raciocínio básico e muito mais \cite{Radford2019}.

Para que o texto possa ser processado numericamente, ele precisa primeiro ser convertido
em tokens e depois em vetores numéricos. Essa conversão ocorre em duas etapas. A primeira
é a tokenização, realizada por algoritmos como o BPE (\textit{Byte Pair Encoding}), que
divide o texto em unidades de comprimento variável de forma a equilibrar cobertura e
eficiência. Por exemplo, a palavra ``inconstitucional'' pode ser dividida em tokens como
``in'', ``constitu'' e ``cional''. O GPT-2 possui um vocabulário de aproximadamente
50.257 tokens. A segunda etapa converte cada token em um vetor numérico denso por meio
de uma camada de \textit{embedding}: o modelo aprende, durante o treinamento, uma
representação vetorial para cada token do vocabulário, de modo que tokens semanticamente
próximos tendam a ter representações próximas no espaço vetorial.

O GPT-2 (\textit{Generative Pre-trained Transformer 2}), proposto por
\citeonline{Radford2019} pela OpenAI em 2019, é um modelo autorregressivo baseado
exclusivamente no componente decodificador do \textit{Transformer} original. Ele foi
pré-treinado em um corpus de aproximadamente 40 GB de texto coletado da internet,
abrangendo uma grande diversidade de domínios e estilos. A variante GPT-2 \textit{small},
utilizada neste trabalho, possui 12 blocos \textit{Transformer} empilhados, 12 cabeças de
atenção por bloco, dimensão de modelo igual a 768 e aproximadamente 124 milhões de
parâmetros treináveis. Apesar de ser a menor variante da família GPT-2, esse número de
parâmetros é cerca de 70 vezes maior do que o de uma ResNet-50 (25 milhões de
parâmetros), o que ilustra a escala característica dos modelos baseados em \textit{Transformer}.

O crescimento em escala dos LLMs ao longo dos anos seguintes foi vertiginoso: o GPT-3,
publicado em 2020, possui 175 bilhões de parâmetros; o PaLM, do Google, chega a 540
bilhões. Mesmo modelos de médio porte como o GPT-2 impõem custos energéticos e
computacionais consideráveis em ambientes de produção, onde a inferência é executada
milhares ou milhões de vezes por dia. A poda estruturada surge, nesse contexto, como uma
estratégia concreta para reduzir esses custos: a remoção de cabeças de atenção
redundantes e de neurônios MLP de baixa importância pode diminuir substancialmente a
contagem de operações por inferência, com impacto controlado sobre a qualidade das
predições \cite{Patterson2021}. 
% Esse potencial de redução do consumo energético via poda
% estruturada em modelos \textit{Transformer} é o foco central dos experimentos conduzidos
% neste trabalho.
